\scanpage{99}
\mspage{1}
\uuline{Théorème}\quad $X/k$, $k$ corps \uuline{alg. clos}

$X$ \deleted{\uuline{\unsure{intègre}}}\inserted{\uline{normal} \uline{irréd}} \uuline{projectif}, $X\subset\mathbf{P}^r_k$

\deleted{$\operatorname{prof}X\geq2$ (\unsure{d'où}} \inserted{$\dim X\geq2$} (d'où $\operatorname{prof}X\geq2$)

$G$ groupe algébrique \uuline{affine} \uuline{sans composantes} $\mathrm{G}_m$.
\inserted{(\unsure{p.ex.} un groupe alg. fini).}

$K=\overset{\deleted{\rule{2.7cm}{.3pt}}}{k(t_1,\ldots,t_r)}$, $\overline{K}$ sa clôture algébrique

$Y_{\overline{K}}$ la section hyperplane générique, définie sur $\overline{K}$.

Alors
\[
\mathrm{H}^1(X,G)\longrightarrow\mathrm{H}^1(Y_{\overline{K}},G)\qquad\uuline{\text{est injectif}}
\]
La structure des groupes algébriques \unsure{comm.} nous ramène aux cas simples $\mathrm{G}_a$, \deleted{$\unsure{\mathrm{G}_m}$}, groupe fini simple, \inserted{$\alpha_p$}, $\mu_p$.

Donc on est ramené à prouver

\uline{Corollaire 1}\quad $\mathrm{H}^1(X,\underline{\mathcal{O}}_X)\longrightarrow\mathrm{H}^1(Y_{\overline{K}},\underline{\mathcal{O}}_{Y_{\overline{K}}})$ est injectif

\margin{N.B. \unsure{Généraliser} au cas d'un morphisme dominant $X\longrightarrow\mathbf{P}^r$}

\deleted{\uline{Corollaire 2}\quad $\mathrm{Pic}(X)\longrightarrow\mathrm{Pic}(Y_{\overline{K}})$ est injectif}

\uline{Corollaire 2}\quad $\pi_1(Y_{\overline{K}})\longrightarrow\pi_1(X)$ est surjectif\quad [C'est \uuline{Bertini}]

\uline{Corollaire 3}\quad $\mathrm{H}^1(X,\alpha_p)\unsure{{-}|\longrightarrow}\mathrm{H}^1(\unsure{X|Y}_{\overline{K}},\alpha_p)$ injectif

\uline{Corollaire 4}\quad $\mathrm{H}^1(X,\mu_p)\unsure{{-}|\longrightarrow}\mathrm{H}^1(Y_{\overline{K}},\mu_p)$ injectif.

On part du corollaire 2, connu grâce à Bertini; \uline{ou} Lefschetz-Grauert [\unsure{utilise} \unsure{déjà que} $\operatorname{prof}X\geq2$]. On en déduit
\[
\mathrm{H}^1(X,\underline{\mathcal{O}}_X)^F\Rightarrow\mathrm{H}^1(X,\mathbf{Z}/p\mathbf{Z})\longrightarrow\mathrm{H}^1(Y_{\overline{K}},\underline{\mathcal{O}}_{Y_{\overline{K}}})^F
\]
est injectif, on est donc ramené pour le cor. 1 de prouver que
\[
\mathrm{H}^1(X,\underline{\mathcal{O}}_X)^{1-F}\unsure{{-}|\longrightarrow}\mathrm{H}^1(Y_{\overline{K}},\underline{\mathcal{O}}_{Y_{\overline{K}}})^{1-F}
\]
est injectif [N.B. en tous cas le noyau est stable sous $F$\ldots] i.e. au corollaire 3. Donc on est ramené \inserted{\uline{En car.} $p>0$} \deleted{\illegible{1}} au cas des cor. 3. et 4. En car $p=0$,
\margin{98}
